% WHERE THE RUN A LOST BREAK MADE REACHES TO.
%
% A word set again without a left boundary begins with its first two letters
% in ONE run, so a break between them cannot be taken and no discretionary is
% drawn for it -- see where-a-rebuilt-run-begins.tex, which settles that much.
% What that file does not ask is where the run those two letters make ENDS,
% and the answer decides everything about the NEXT break in the word:
%
%   - the run it makes begins where that one left off, at the letter after
%     the lost break, and not at its own letter;
%   - so it replaces only what stands from there, not the whole rest of the
%     word;
%   - and the letter in front of it is already laid down, so what goes before
%     the break is the hyphen ON ITS OWN.
%
% Measured against pdfTeX with trip's font, whose ligature program puts a
% whole run of characters between every pair of letters. `\noboundary111'
% with breaks wanted at both letters draws one discretionary, and it is the
% SECOND break's:
%
%   ref:    \discretionary replacing 2
%           .\rip -
%           |\rip 5 (ligature |1)
%           |\rip 7 (ligature |)
%
% This engine drew `replacing 67' with `1', `1 (ligature -)', `7 (ligature |)'
% in front of the break, because it began the run at the lost break's own
% letter -- sixty-five nodes early. That is what trip's first pass was
% weighing when it made the break after `[]11-1' come out at b=2 where the
% reference has b=42: the line was a pointful too wide, so it had to shrink
% where the reference had to stretch, and decent came out where loose
% belonged. See what-a-break-is-weighed-as-in-trips-first-pass.tex.
%
% [A] is the case itself. [B] is the same word WITH the boundary, where no
% break is lost and both discretionaries stand. [C] is four letters, where
% the run reaches past two lost breaks.
\catcode`\{=1 \catcode`\}=2
\tracingonline=1 \showboxbreadth=9999 \showboxdepth=9999
\font\rip=trip \rip
\hsize=1pt \parindent=0pt \pretolerance=-1 \tolerance=10000 \hbadness=10000
\hyphenchar\rip=`-
\lccode`1=`1 \patterns{111}
\lefthyphenmin=1 \righthyphenmin=1
\message{[A a lost break, and where the next one begins]}
\setbox1=\vbox{\noindent\ \noboundary111 \par}\showbox1
\message{[B the same word with the boundary, where nothing is lost]}
\setbox1=\vbox{\noindent\ 111 \par}\showbox1
\message{[C four letters, without the boundary]}
\setbox1=\vbox{\noindent\ \noboundary1111 \par}\showbox1
\message{[done]}
\end
